% TOP_PROBLEM: {"id": 1100110, "title": "Questions in Geometric Group Theory — Q 1.10", "category_id": 17, "category": {"id": 17, "name": "group_theory", "display_name": "Group Theory", "description": "Problems about groups, group actions, representations, and related algebraic structures.", "slug": "group-theory", "order_index": 17, "created_at": "2026-07-31T15:26:25.671Z"}, "status": "partially_solved", "problem_number": "AMR-010-0110", "set_id": 13}
% FIRST_POSTED: 2026-10-01
% ENTRY_KIND: statement_only
% UnsolvedMath ID 1100110; source catalogue code AMR-010-0110
% !TeX program = lualatex
% Definition and sources only; this file is not an LLM solution attempt.
\documentclass[11pt,a4paper]{article}
\usepackage[margin=25mm]{geometry}
\usepackage{fontspec}
\setmainfont{lmroman10-regular.otf}[BoldFont=lmroman10-bold.otf,
 ItalicFont=lmroman10-italic.otf,BoldItalicFont=lmroman10-bolditalic.otf]
\usepackage{amsmath,amssymb,mathtools}
\usepackage{unicode-math}
\setmathfont{latinmodern-math.otf}
\AtBeginDocument{\let\setminus\smallsetminus}
\usepackage{xurl,hyperref}
\hypersetup{colorlinks=true,urlcolor=blue}
\setcounter{secnumdepth}{0}
\setlength{\parindent}{0pt}
\setlength{\parskip}{5pt}
\setlength{\emergencystretch}{3em}
\begin{document}
\section*{Questions in Geometric Group Theory — Q 1.10}\label{1100110}
{\small UnsolvedMath ID 1100110 · AMR-010-0110 · Group Theory · AMR Open Problem Lists}
\subsection{Definitions and mathematical statement}
Let $G$ be a finitely generated word-hyperbolic group. Thus, for a finite symmetric generating set, its Cayley graph with unit edges has $\delta$-thin geodesic triangles for some finite $\delta$. Let $H\leq G$ be a finitely presented subgroup, with its own finite generating set and finite list of defining relations.

Assume that every element of $G$ has a positive power lying in $H$:
\[
 \forall g\in G\ \exists n(g)\in\mathbb Z_{\geq1}
          \quad g^{n(g)}\in H.
\]
Must the index $[G:H]$, the number of left cosets of $H$, be finite?

The exponent is allowed to depend on $g$; there is no common exponent in the hypothesis. The subgroup is not assumed normal. Consequently the assertion cannot in general be rephrased as a torsion statement about a quotient group $G/H$, since that quotient need not be a group. When $H$ is normal, the displayed condition does say that every element of the quotient has finite order, but finite presentation of the subgroup is still a separate hypothesis.

The reverse implication is elementary: if $H$ has finite index, the sequence of cosets $H,gH,g^2H,\ldots$ repeats, forcing a positive power of $g$ into $H$. The question asks for the converse in the stated hyperbolic and finite-presentation setting. It does not merely ask whether $H$ meets every conjugacy class or contains powers of a fixed generating set.
\subsection{Short English statement}
Must a finitely presented subgroup of a hyperbolic group have finite index if it contains a positive power of every group element?
\subsection{Sources}
\begin{itemize}
\item[S1] ULAM.AI and the UnsolvedMath Contributors, supplied problems.json, record 1100110 (AMR-010-0110), AMR Open Problem Lists. Catalogue identity and source transcription; CC BY 4.0.
\url{https://huggingface.co/datasets/ulamai/UnsolvedMath}.
\item[S2] Bestvina - Questions in Geometric Group Theory (pdf) (2004), Question 1.10 (PDF page 3). Original source indexed by AMR-010-0110.
\url{https://www.math.utah.edu/~bestvina/eprints/questions-updated.pdf}.
\end{itemize}
\end{document}
